\clearpage\begingroup\NoSourceHead\SourceLeaf{638}
\HistoricalFoot{}{1\textsuperscript{*}}
\makeatletter\tagsleft@true\makeatother
\renewcommand{\thefootnote}{\arabic{footnote})}\setcounter{footnote}{0}
\begin{center}{\bfseries I. Nachweis der Erdrotation.}\\[5pt]
\textls[100]{1. Fallversuche und ihre Theorie vor \textsc{Gauss}.}\end{center}
\noindent\hspace*{\parindent}\textsc{Newton} war der erste, der den Gedanken aussprach, man könne die Erdrotation
experimentell dadurch nachweisen, dass man Kugeln von einer beträchtlichen
Höhe herabfallen lässt, die dann um ein Weniges östlich von
der Lotlinie des Ausgangspunktes auf der Erde ankommen müssen; \textsc{Newton}
teilte diesen Gedanken im Nov. 1679 \textsc{Hooke}, dem damaligen Sekretär der
Royal Society of London, mit und schlug vor, Versuche in dieser Richtung
durchzuführen. Der Gedanke ward von der Akademie zwar mit grossem Beifall
aufgegriffen, aber \textsc{Hooke} antwortete zunächst ausweichend, und als er
gedrängt wurde, seiner Pflicht nachzukommen, stellte er seine Versuche nur
bei einer Fallhöhe von 27 Fuss an, bei der selbst für die heutigen Mittel
die östliche Abweichung nicht zu konstatieren ist; in den Akademieberichten
findet sich nichts weiter von derartigen Versuchen, woraus man schliessen
muss, dass sie negativ ausfielen. Obgleich dieser \textsc{Newton}sche Gedanke auch
in der Folgezeit lebendig blieb, waren die experimentellen Schwierigkeiten
so gross, dass erst 113 Jahre später (1792) der junge Abbé \textsc{Guglielmini} in
Bologna es unternahm, Fallversuche in der Torre degli Asinelli anzustellen\footnote{\textsc{Joh. Baptista Guglielmini}, \textit{De diurno terrae motu experimentis physico-mathematicis confirmato opusculum}. Bononiae 1792. Das Buch erschien mit Approbation der obersten Kirchenbehörde, zeigt also,
dass diese schon damals die Erdrotation als mit dem Dogma vereinbar ansah.}.
\textsc{Guglielmini} hatte damals noch mit grossen Schwierigkeiten zu kämpfen; einerseits
weist der verwandte Turm viele Löcher auf, durch die ein die fallenden Kugeln
ablenkender Luftzug entstand, zweitens müssen die Kugeln vor dem Ablassen
vollkommen ruhig hängen, und dürfen beim Loslassen keine Geschwindigkeit
mitbekommen, was \textsc{Guglielmini} durch eine zangenartige Vorrichtung erreichte,
bei deren Öffnen der die Kugel haltende Faden losgelassen wurde; schliesslich
\clearpage\SourceLeaf{639}\ArithmeticHead{4}{L}{H. GEPPERT, ÜBER GAUSS’ ARBEITEN ZUR MECHANIK UND POTENTIALTHEORIE.}
\noindent mass \textsc{Guglielmini} nicht die Höhe, sondern die Zeit des Falles, deren Bestimmung
damals noch besonders schwierig war. Schwere Metallkugeln durchfielen
eine Strecke von 240 par. Fuss $=77{,}96$ m und schlugen unten ihre
Eindrücke in eine Wachsplatte; es wurden bei Nacht und ruhiger Luft 16
Versuche angestellt, aus denen \textsc{Guglielmini} die östliche Abweichung des Einschlagspunktes
vom scheinbaren Lote zu 8,375 Pariser Linien $=18{,}89\pm2{,}5$ mm
(bei Weglassung eines fehlerhaften Versuches ergibt sich 16,69 mm) bestimmte;
ausserdem fand sich jedoch eine südliche Abweichung von 5,272 Linien
$=11{,}89\pm1{,}1$ mm. Leider sind diese mit soviel Mühe angestellten Versuche
nicht beweiskräftig, denn sie enthalten, wie \textsc{Benzenberg} zuerst bemerkte, einen
systematischen Fehler; \textsc{Guglielmini} bestimmte nämlich ungünstiger Witterungsverhältnisse
wegen die Lage des scheinbaren Lotes erst ein halbes Jahr nach
seinen Fallversuchen, und es ist die Annahme jedenfalls nicht von der Hand
zu weisen, dass der an sich schiefe Turm in Folge der veränderten Temperatur
und Sonnenstrahlung in dieser Zeit eine Veränderung der Krümmung
erfahren habe. Trotz der grossen Vorsichtsmassregeln sind die \textsc{Guglielmini}schen
Versuche unbefriedigend, da der Unterschied zwischen den extremen Abweichungen
sich auf die Grösse von ca. 14 mm beläuft.

\textsc{Guglielmini} hatte seine Resultate u. a. an \textsc{Laplace} geschickt, und dies
mag letzterem Veranlassung gewesen sein, die theoretische Seite dieser Fallversuche
zu untersuchen. Er entwickelte im Jahre 1803 die Theorie der
Bewegung eines fallenden Körpers auf der rotierenden Erde\footnote{\textsc{Laplace}, \textit{Mémoire sur le mouvement d’un corps, qui tombe d’une grande hauteur}. Bull. de
Science par la Société Philomath., Prairial an 11 (1803).}, und seine
schwer zugängliche Abhandlung wurde zwei Jahre später in etwas veränderter
Fassung in seine Mécanique céleste aufgenommen\footnote{\textsc{Laplace}, \textit{Mécanique céleste}, 4, livre X, chap. 5, 1805, vergl. Oeuvres 4, 1880, S. 295—306. Auch
abgedruckt in \textsc{Benzenberg} a. a. O. \textsuperscript{5)} S. 388—400.}. \textsc{Laplace} besass diese Theorie
schon etwas früher und teilte ihre Resultate \textsc{Guglielmini} mit, der seinerseits
eine der später zu besprechenden \textsc{Olbers}schen ähnliche Überlegung entwickelt
hatte. Werfen wir einen kurzen Blick auf die \textsc{Laplace}sche Theorie. Der
Ausgangspunkt der Fallbewegung wird auf zwei verschiedene Koordinatensysteme
bezogen, erstens auf ein mit der Erde mitrotierendes rechtwinkliges System
$x,y,z$, dessen Anfangspunkt mit dem Erdmittelpunkt, dessen $x$-Achse mit
\clearpage\SourceLeaf{640}\ArithmeticHead{5}{R}{I. NACHWEIS DER ERDROTATION.}
\noindent der Rotationsachse der Erde zusammenfällt, zweitens auf ein festes Polarkoordinatensystem
$r,\theta,\varpi$, worin $\theta$ das Komplement der geographischen Breite bedeutet.
Die Anwendung des \textsc{d’Alembert}schen Prinzips ergibt dann drei Differentialgleichungen
2. Ordnung in $r,\theta,\varpi$. Während des Falles mögen diese
drei Grössen die Veränderungen $\delta r,\,\delta\theta,\,\delta\varpi+nt$ erleiden, worin $n$ die Winkelgeschwindigkeit
der Erdrotation und $t$ die Fallzeit bezeichnet. Da $n$ sehr
klein ist, integriert \textsc{Laplace} nicht die exakten, sondern die bis auf Grössen
2. Ordnung gekürzten Bewegungsgleichungen:
\[
\begin{gathered}[t]
\frac{d^2\delta r}{dt^2}+S\frac{d\delta r}{dt}+g=0,\quad
\frac{d^2\delta\theta}{dt^2}+S\frac{d\delta\theta}{dt}-g\frac{\partial y}{d\theta}=0,\\[5pt]
\frac{d^2\delta\varpi}{dt^2}\sin\theta+2n\sin\theta\frac{d\delta r}{dt}
+S\sin\theta\frac{d\delta\varpi}{dt}-\frac{g}{\sin\theta}\frac{\partial y}{\partial\varpi}=0.
\end{gathered}\tag{1}
\]
\noindent Hierin ist auch die Luftreibung in Ansatz gebracht, indem $S$ eine von $r,\delta r$
abhängende Reibungsfunktion bedeutet. Unter Berücksichtigung der Anfangsbedingungen
ergibt die zweite Gleichung sofort die Lösung
\[
\delta\theta=-\delta r\frac{\partial y}{\partial\theta},\tag{2}
\]
\noindent was physikalisch bedeutet, dass der fallende Körper von dem scheinbaren Lote
keine merkliche Abweichung nach Süden aufweist. Setzt man die Reibung
proportional dem Quadrat der Geschwindigkeit, und bedeutet $m$ den entsprechenden
Reibungsfaktor, so lässt sich die erste Gleichung (1) exakt in
geschlossener Form integrieren; die Reihenentwicklung der Lösung lautet
\[
\delta r=-\frac{g}{2}t^2+m\frac{g^2}{12}t^4-m^2\frac{g^3}{45}t^6+\mathinner{\cdotp\cdotp\cdotp}.\tag{3}
\]
\noindent Damit ergibt schliesslich die Integration der dritten Gleichung (1) die Abweichung
des fallenden Körpers vom scheinbaren Lote in der Ostrichtung zu:
\[
\tfrac{1}{3}ngt^3\sin\theta\left(1-\frac{mg}{4}t^2+\frac{61}{840}m^2g^2t^4-\mathinner{\cdotp\cdotp\cdotp}\right).\tag{4}
\]
\noindent Ist $h$ die Fallhöhe, so beträgt die östliche Abweichung des Aufschlagspunktes
vom scheinbaren Lote des Ausgangspunktes
\[
\tfrac{2}{3}nh\sqrt{\frac{2h}{g}}\sin\theta,\tag{5}
\]
\noindent unter Vernachlässigung der quadratischen Glieder von $m$ und $n$. Dagegen
ist die südliche Abweichung des Aufschlagspunktes unmerklich.
\clearpage\SourceLeaf{641}\ArithmeticHead{6}{L}{H. GEPPERT, ÜBER GAUSS’ ARBEITEN ZUR MECHANIK UND POTENTIALTHEORIE.}
\noindent\hspace*{\parindent}Wendet man die \textsc{Laplace}sche Formel (5) auf die \textsc{Guglielmini}schen Resultate
an, so ergibt sich der theoretische Wert der östlichen Abweichung zu
10,78 mm, sodass also der experimentelle Wert gegenüber dem theoretischen
zu gross ausfällt; die von \textsc{Guglielmini} gefundene südliche Abweichung ist mit
der Theorie unvereinbar\footnote{\textsc{Laplace} führt in seiner Note weiterhin die Theorie einer von der Erde aus senkrecht nach oben
geschossenen Kanonenkugel durch und zeigt, dass diese mit einer westlichen Abweichung vom Ausgangspunkte
einschlägt, wohingegen auch hier die südliche Abweichung verschwindend gering ist. Wegen experimenteller
Schwierigkeiten ist dieses Ergebnis praktisch nie bestätigt worden. Auch Fallversuche wurden
in Paris nicht unternommen.}.

Die Dissonanz zwischen Experiment und Theorie veranlasste \textsc{Benzenberg}
im Jahre 1802, die Fallversuche wieder aufzunehmen\footnote{\textsc{J. Fr. Benzenberg}, \textit{Versuche über die Umdrehung der Erde}, Dortmund 1804.}. Die Versuche wurden
im grossen Michaelisturm in Hamburg bei einer Fallhöhe von 235 Fuss $=$
76,34 m angestellt; dieser Turm eignete sich zu den Fallversuchen insbesondere
deshalb, weil er eine grosse Fallhöhe gestattet und überdies geschlossen ist,
sodass störender Luftzug verhältnismässig abgeschwächt ist. \textsc{Benzenberg} verwendet
wesentlich die \textsc{Guglielmini}sche Apparatur und ist bestrebt, die Fehlerquellen
seines Vorgängers sorgfältig zu vermeiden. Die verwendeten kleinen
Bleikugeln schlugen ihre Eindrücke in ein mit Kreide überzogenes Hartholzbrett.
Es wurden 16 brauchbare Versuche angestellt, aus denen sich die
östliche Abweichung zu $9\pm3{,}6$ mm und eine südliche Abweichung von
$3{,}4\pm2{,}5$ mm errechnete. Trotz aller Mühe sind die \textsc{Benzenberg}schen Versuche
schlechter als die seines Vorgängers und kaum beweiskräftig; schwanken
doch die Abweichungen zwischen 47 mm nach Osten und 31,5 mm nach
Westen, und von den Abweichungen im Meridian fielen $\frac{2}{5}$ nach Norden;
auch gestatteten die Verhältnisse nicht einmal die Anwendung eines Mikroskops,
um die absolute Ruhe im Ausgangspunkte sicherzustellen. Die von
\textsc{Benzenberg} bestimmte östliche Abweichung harmoniert gut mit dem von der
\textsc{Laplace}schen Formel (5) geforderten Wert von 8,7 mm, dagegen steht die
südliche Abweichung wie bei \textsc{Guglielmini} im Widerspruch zur Theorie.

\textsc{Benzenberg} sandte seine Versuche im Herbst 1802 an \textsc{Olbers} nach Bremen,
damit er sie mit der Theorie vergliche. \textsc{Olbers} kannte zwar die \textsc{Laplace}sche
Abhandlung dem Namen nach, muss sie aber nicht gegenwärtig gehabt
haben, denn er entwickelte von sich aus eine neue Theorie, die andere Werte
\clearpage\SourceLeaf{642}\ArithmeticHead{7}{R}{I. NACHWEIS DER ERDROTATION.}
\noindent als die \textsc{Laplace}sche liefert und sich später als falsch erwies. Er errechnete
für die \textsc{Benzenberg}schen Versuche eine östliche Abweichung von 13,5 mm
und eine südliche von 3,5 mm und geriet damit in Unstimmigkeit einerseits
gegenüber den Versuchsresultaten, andererseits gegenüber den \textsc{Laplace}schen
Ergebnissen; dies war für \textsc{Olbers} die Veranlassung, die ganze Frage in einem
Briefe vom 12. November 1802 \textsc{Gauss} vorzulegen\footnote{Briefwechsel zwischen \textsc{Olbers} und \textsc{Gauss}. \textit{Wilh. Olbers Leben und Werke} II, 1, Berlin 1900,
S. 107 ff.}. \textsc{Olbers} behauptete darin
eine südliche Abweichung jedes fallenden Körpers in Folge des Widerstandes
der Luft, und bestimmte deren Grösse zu
\[
\frac{n^2}{4}r\sin2\varphi\left(t^2-\frac{2h}{g}\right),\tag{6}
\]
\noindent worin $r$ den Erdradius, $\varphi$ die geographische Breite und $t$ die Fallzeit bedeutet.
Da im widerstehenden Mittel $t^2>\frac{2h}{g}$ ist, ergab sich nach \textsc{Olbers}
eine durch den Widerstand erzeugte südliche Abweichung der Grössenordnung
$n$, also durchaus beobachtbar. \textsc{Gauss} widersprach zwar in seinem Briefe vom
18. Nov. 1802 sofort den \textsc{Olbers}schen Überlegungen, aber seine flüchtigen
Einwürfe waren nicht stichhaltig. Erst ein Brief, den \textsc{Benzenberg} kurz darauf
an \textsc{Gauss} richtete, bewog letzteren, sich mit der Frage gründlich auseinanderzusetzen
und die Fundamentalgleichungen für die Bewegung schwerer
Körper auf der rotierenden Erde neu zu entwickeln; \textsc{Gauss} fasste seine Ergebnisse
in einer kleinen Abhandlung zusammen, die in \textsc{Benzenberg}s Buche
(S. 363—371) veröffentlicht wurde\footnote{Abgedruckt Werke V, S. 498—503.}. Im Briefe vom 1. März 1803 teilt \textsc{Gauss}
seine endgültigen Resultate \textsc{Olbers} mit.
\clearpage\endgroup

